% ============================================================
%   Chapter 7  \textemdash\  Conclusions and Future Work
% ============================================================

\chapter{Conclusions and Future Work}\label{ch:conclusions}

This chapter closes the document. Section~\ref{sec:concl-contributions} recapitulates the
contributions, Section~\ref{sec:concl-objectives} traces them against the specific objectives
laid out in Chapter~\ref{ch:introduction}, Section~\ref{sec:concl-lessons} reflects on the
lessons learned during the project, and Section~\ref{sec:concl-future} outlines research
directions that go beyond the scope of a Bachelor's thesis.

\section{Summary of contributions}\label{sec:concl-contributions}

The work presented in this document produced the following concrete contributions:

\begin{enumerate}
  \item \textbf{CALA-SafeRL framework.} A modular Safe RL framework for autonomous driving on
        CARLA 0.9.16, organised as a chain of interchangeable \texttt{gymnasium.Wrapper} layers.
        The wrapper chain isolates the PPO agent, the safety shield, the reward shaper and the
        simulator from each other, which makes each component independently testable and
        replaceable.
  \item \textbf{Semantic 739-dimensional observation space.} A sensor pre-processing stage
        packages three LIDAR channels (combined, dynamic-only, static-only) of \(240\) rays each
        together with \(19\) geometric and route features (including lane-marking type flags)
        derived from CARLA's Waypoint API. The decomposition of the LIDAR into semantic channels
        mitigates the zig-zag pathology introduced by static guardrails.
  \item \textbf{PPO agent with modern stabilisers.} An implementation of PPO with Generalised
        Advantage Estimation, clipped surrogate updates, value clipping, approximate-KL early
        stopping, cosine learning-rate annealing, return normalisation and a Welford running-mean
        estimator that is applied only to the vector block of the observation.
  \item \textbf{Basic safety shield.} A lightweight, interpretable shield that combines LIDAR
        thresholding with lane-keeping corrections derived from the Waypoint API. The shield
        obeys the minimal-interference principle: its override action is zero whenever the
        situation is safe and grows smoothly with the severity of the violation.
  \item \textbf{Adaptive-horizon safety shield.} A richer shield that classifies the current
        situation into three risk levels, predicts the vehicle's future trajectory with a
        kinematic bicycle model calibrated for the Tesla Model 3, and selects the least invasive
        corrective action from a priority-ordered candidate set. The shield supports a
        pedestrian emergency, a lane-change relaxation rule and a threshold-floor safeguard that
        prevents the \emph{critical} level from rejecting all candidate actions.
  \item \textbf{Dense reward shaper.} A reward wrapper with twelve additive components plus an
        alive bonus, all of them gated by speed when appropriate to avoid the degenerate
        ``stop-and-collect'' optimum. The shaper also implements a lane-change transition guard
        and a shield grace period that stabilise the interaction between the agent, the shield
        and the reward.
  \item \textbf{Performance-driven curriculum with rollback.} A curriculum manager that divides
        the interval \([0, N_{\max}]\) of NPC vehicles into five uniform stages and advances or
        rolls back according to \(100\)-episode rolling-window statistics. Stage transitions
        update the NPC count in place, without reconnecting to CARLA.
  \item \textbf{Live metrics pipeline.} A SQLite-based logger in WAL mode that exposes two
        tables (\emph{per-episode} and \emph{per-update}) and a Streamlit dashboard that reads
        them while training is still in progress. Unknown metric keys are rendered as
        first-class citizens, which makes the dashboard forward-compatible with future
        experiment extensions.
  \item \textbf{Reproducibility and tooling.} Unit tests for the reward shaper that do not
        require a running simulator, a \texttt{ruff}-enforced code style, seed-separated
        training and evaluation protocols, a checkpoint format that includes the observation
        normaliser state, and a reproducibility guide in Appendix~\ref{app:reproducibility}.
\end{enumerate}

\section{Fulfilment of objectives}\label{sec:concl-objectives}

Table~\ref{tab:objectives-traceability} traces the specific objectives of
Section~\ref{sec:intro-objectives} against the chapters in which each of them is addressed and
against the contributions enumerated above.

\begin{table}[h]
  \centering
  \caption{Traceability matrix between specific objectives and the rest of the document.}
  \label{tab:objectives-traceability}
  \small
  \begin{tabular}{@{}llll@{}}
    \toprule
    \textbf{Objective} & \textbf{Chapter(s)} & \textbf{Contribution} & \textbf{Status} \\
    \midrule
    SO1 \textendash\ Gymnasium CARLA wrapper
         & \ref{ch:design}, \ref{ch:implementation}
         & 1, 2
         & Achieved. \\
    SO2 \textendash\ PPO agent with stabilisers
         & \ref{ch:design}, \ref{ch:implementation}
         & 3
         & Achieved. \\
    SO3 \textendash\ Basic shield
         & \ref{ch:design}, \ref{ch:implementation}, \ref{ch:experimentation}
         & 4
         & Achieved. \\
    SO4 \textendash\ Adaptive-horizon shield
         & \ref{ch:design}, \ref{ch:implementation}, \ref{ch:experimentation}
         & 5
         & Achieved. \\
    SO5 \textendash\ Dense reward shaper
         & \ref{ch:design}, \ref{ch:implementation}
         & 6
         & Achieved. \\
    SO6 \textendash\ Curriculum with rollback
         & \ref{ch:design}, \ref{ch:implementation}, \ref{ch:experimentation}
         & 7
         & Achieved. \\
    SO7 \textendash\ Live metrics pipeline
         & \ref{ch:design}, \ref{ch:implementation}
         & 8
         & Achieved. \\
    \bottomrule
  \end{tabular}
\end{table}

The quantitative evaluation that anchors \emph{``Achieved''} for SO3, SO4 and SO6 is deferred to
the final reports of Chapter~\ref{ch:experimentation}; until the TBD entries of that chapter
are filled in, the empirical side of the argument rests on the acceptance criteria of
Table~\ref{tab:acceptance-criteria}.

\section{Lessons learned}\label{sec:concl-lessons}

The development of the framework surfaced a number of design lessons that, in retrospect, would
have saved considerable time if they had been internalised earlier in the process.

\begin{description}
  \item[\textbf{Place the shield above the reward shaper.}] An earlier prototype inverted this
        ordering and the shaper could not read the executed action. The reward-shaping signals
        that depend on the shield (intervention penalty, smoothness suppression during
        override, grace period) were all computed incorrectly until the chain was reordered.
  \item[\textbf{Never normalise the LIDAR.}] Applying the running-mean normaliser to the
        \(720\) LIDAR dimensions was the single biggest source of training instability.
        When the curriculum advances from stage~0 (no NPCs, the dynamic channel is constant at
        \(1.0\)) to stage~1, the running standard deviation of the dynamic channel is close to
        zero and the normalised inputs become extreme outliers. Restricting the normaliser to
        the \(15\)-dimensional vector block eliminated the gradient spikes without any further
        change to the training loop.
  \item[\textbf{Gate the alive bonus by speed.}] An unconditional alive bonus, even a small one,
        opens a local optimum in which the agent stops, does nothing, and earns
        \(0.15\)/step forever. Gating the alive bonus by the speed factor \(g_t\) of
        Equation~\eqref{eq:speed-gate} closes this trap without hurting exploration.
  \item[\textbf{Give the agent a grace period after a shield override.}] A brake-based shield
        intervention drops the vehicle's speed, which would immediately trigger the idle penalty
        and push the agent into a stationary equilibrium. The shield-grace mechanism of
        Section~\ref{sec:design-reward} breaks this feedback loop. Critically, the grace counter
        is re-armed only when the shield fires \emph{and} the counter is already at zero;
        otherwise, an always-active shield would perpetually reset the grace period and silence
        the idle penalty forever.
  \item[\textbf{Prefer in-place updates to reconnection.}] Every early attempt at changing the
        curriculum stage by rebuilding the environment caused CARLA to deadlock. Updating the
        attribute \texttt{num\_npc\_vehicles} on the base environment and letting the change
        take effect at the next \texttt{reset} is both simpler and robust.
  \item[\textbf{Attribute credit to the executed action, not to the proposed one.}] Storing the
        action that the shield let through, and recomputing its log probability under the
        current policy, is the single correction that made shielded PPO actually learn.
        Without it, the gradient estimator pushes the policy towards an action that never
        reached the environment, and training stagnates.
\end{description}

\section{Future work}\label{sec:concl-future}

Several extensions of this framework are natural next steps. They are listed below in order of
increasing research effort.

\begin{itemize}
  \item \textbf{Dynamic bicycle model.} The kinematic model used by the adaptive shield ignores
        tyre slip and lateral dynamics. At higher speeds or in tight turns this approximation
        becomes inaccurate and the predicted trajectory can diverge from what the simulator
        actually produces. Replacing it with a dynamic bicycle model or with a simple
        lateral-dynamics model~\cite{rajamani2012vehicle} would tighten the safety guarantees
        without changing the rest of the architecture.
  \item \textbf{Reachability-based shield.} A principled alternative to the candidate-search
        strategy is Hamilton\textendash Jacobi reachability
        analysis~\cite{bansal2017hj,fisac2019safety}, which precomputes the set of unsafe
        states and projects the agent's action onto the largest safe control invariant set.
        This yields formal safety guarantees but imposes a non-trivial offline computation cost.
  \item \textbf{Camera-based perception.} The observation space currently does not include RGB
        imagery or semantic-segmentation maps. Adding a convolutional perception head trained
        jointly with the policy, or reusing the frozen backbone of an imitation-learning
        pipeline~\cite{chitta2023transfuser}, would enable the agent to perceive traffic lights,
        road signs and crosswalks that are not visible from the semantic LIDAR alone.
  \item \textbf{Urban benchmarks.} The evaluation has been concentrated on Town04, a highway
        map. Urban maps such as Town01\textendash Town03 and Town10HD feature intersections,
        traffic lights and dense pedestrian traffic that would stress-test the pedestrian-
        emergency branch of the adaptive shield. Adapting the curriculum to this richer
        scenario is itself a research question.
  \item \textbf{Multi-agent safe RL.} The current setup has a single RL-controlled ego vehicle
        surrounded by rule-based NPCs. Replacing the NPCs with independent RL agents and
        composing their individual shields into a global safety filter is an open research
        direction at the intersection of Safe RL and multi-agent RL.
  \item \textbf{Constrained policy optimisation.} PPO is used here as a black-box optimiser of
        the shaped reward. A closer integration with Safe RL would replace the reward shaping
        by explicit safety constraints and solve the resulting Constrained MDP with CPO or a
        Lagrangian variant~\cite{achiam2017cpo,stooke2020responsive}. This would remove the
        hand-tuning of reward weights at the cost of a more complex optimiser.
  \item \textbf{Benchmarking against other shields.} The MetaDrive and HighwayEnv ecosystems
        offer a variety of shielding baselines. Porting those baselines to CARLA and comparing
        them on a shared benchmark would provide an external reference for the adaptive shield
        proposed here.
\end{itemize}

The framework has been designed with these extensions in mind. Each of them should be achievable
by introducing a new wrapper in the chain of Figure~\ref{fig:wrapper-chain} or by replacing the
existing \texttt{Shield} or \texttt{RewardShaper} classes, without modifying the remaining
components.
